% ==========================================================================
%  Section 7 -- Finite Fourier marginals.  Owner: W07.  Target 1100 words.
%  Frozen labels (CONTRACT section 4): sec:marginals, def:cylinder-map,
%  thm:char-function, thm:finite-family, lem:hC-identification.  Added here:
%  the eq: labels and rem:index-dictionary.
%
%  What this file does.  Section 4 proves a limit of PAIRINGS,
%  \mathcal Q_N(probe,probe) -> Real Q(probe,probe); the sources then assert that the
%  covariance MATRIX of the paired cosine/sine family is the resolvent
%  matrix without deriving it (ledger GAP R2).  That passage is derived here,
%  entry by entry, in lem:hC-identification: no new limit theorem is needed,
%  because thm:covariance-limit already identifies the limit of exactly these
%  pairings.  The section keeps apart the two things a reader tends to merge,
%  the finite-N matrix (built from Q_N) and the target matrix (built from Q).
%
%  Evidence: everything is (M), except two standard facts named inside the
%  proof of thm:finite-family.  No (F) tag appears: the finite-family
%  interface declarations are not in the generated ledger, so no \lean{}
%  name may be printed here (audit check N).
% ==========================================================================
\section{Finite Fourier marginals}\label{sec:marginals}

This section takes the covariance limit of \cref{sec:covariance} from pairings
to matrices. \Cref{def:cylinder-map} fixes the finite coordinate maps and
records that they are continuous; \cref{thm:char-function} observes that every
finite coordinate law of the lattice sequence is centred Gaussian, and is
therefore nothing but its covariance matrix; \cref{lem:hC-identification}
computes that matrix from the Green pairing of \cref{thm:covariance-limit}; and
\cref{thm:finite-family} concludes convergence of the finite-dimensional laws
for every finite family. The order of the argument is what needs care, and
\cref{lem:hC-identification} is where it sits.

\subsection{Coordinates on the carrier}

\begin{definition}[Cylinder maps]\label{def:cylinder-map}
Let $r\ge1$. A \emph{finite family} is a tuple
$F=((k_1,\varepsilon_1),\dots,(k_r,\varepsilon_r))$ with $k_j\in\Z^2$ and a
label $\varepsilon_j\in\{\mathrm c,\mathrm s\}$ at each slot; \emph{both labels
are allowed at every mode, the zero mode included}. Put $\psi_j:=\cos_{k_j}$ if
$\varepsilon_j=\mathrm c$ and $\psi_j:=\sin_{k_j}$ if
$\varepsilon_j=\mathrm s$, with the cosine and sine tests of
\cref{sec:tightness}. The \emph{cylinder map} of $F$ is
$\Pi_F:\mathcal C\to\R^r$,
\begin{equation}\label{eq:cylinder-map}
  (\Pi_Fc)_j:=
  \begin{cases}
    \Re c_{k_j}, & \varepsilon_j=\mathrm c,\\[2pt]
    -\Im c_{k_j}, & \varepsilon_j=\mathrm s .
  \end{cases}
\end{equation}
The sign in the second line is the one carried by the phase of the lift
\eqref{eq:lift}, and it is bookkeeping: for every configuration $\omega$ and
every slot $j$,
\begin{equation}\label{eq:cylinder-lift}
  \bigl(\Pi_F(L_N\omega)\bigr)_j=\sqrt{w_3(k_j)}\;\omega(\psi_j),
\end{equation}
because that phase is $e^{-2\pi i\,k\cdot(ax)/L}$, whose real part is the
cosine test and whose negative imaginary part is the sine test. So a cylinder
coordinate is $\sqrt{w_3(k_j)}$ times the value of the corresponding real
test, and $R_F:=\operatorname{diag}\bigl(\sqrt{w_3(k_1)},\dots,
\sqrt{w_3(k_r)}\bigr)$ is the fixed diagonal scaling from real-test
coordinates to cylinder coordinates.

Each coordinate in \eqref{eq:cylinder-map} is a scalar multiple of one of the
coordinate functionals of \cref{lem:coordinate-embedding}, hence continuous,
and $\Pi_F$ is continuous into $\R^r$. A continuous map is Borel, so for a
Borel probability measure $\mu$ on $\mathcal C$ the image
$\mu^F:=\mu\circ\Pi_F^{-1}$ is a Borel probability measure on $\R^r$, the
\emph{$F$-marginal} of $\mu$. The family carrying both labels at each of the
modes $k_1,\dots,k_r$ is the complex tuple $(c_{k_1},\dots,c_{k_r})$ read slot
by slot.

The sine test at the zero mode is identically zero, and its slot is retained
deliberately rather than discarded as empty: the coordinate it records is
$-\Im c_0$, and
\eqref{eq:cylinder-map} shows that no other slot sees that coordinate. The
family is therefore a determining class, which is what
\cref{thm:cylinder-uniqueness} consumes; dropping the slot would leave the
cylinder sigma algebra strictly smaller than $\mathcal B(\mathcal C)$ and the
characterisation of \cref{thm:main} would fail. The slot costs nothing
elsewhere: \eqref{eq:entries} gives it covariance entry $0$ and
\cref{thm:char-function} therefore gives it the degenerate law $\delta_0$,
which is the truth, since $(\mathcal L_N\omega)(0)=a^{2}\sum_{x}\omega(x)$ is
real for every configuration.
\end{definition}

\subsection{The finite-mode laws and their covariance matrix}

\begin{theorem}[Characteristic functions of the finite-mode laws]\label{thm:char-function}
Let $F$ be a finite family with $r$ slots and let $\mu_N$ be the law
\eqref{eq:finite-source-law}. Then:
\begin{enumerate}
\item[(i)] for every $N$ the marginal $\mu_N^F$ is a centred Gaussian measure on
  $\R^r$, possibly degenerate, with covariance matrix
  \begin{equation}\label{eq:sigma-N}
    \Sigma_{F,N}(j,j'):=\sqrt{w_3(k_j)w_3(k_{j'})}\;\Real \mathcal Q_N(\psi_j,\psi_{j'}),
  \end{equation}
  and hence with characteristic function
  $t\mapsto\exp\bigl(-\tfrac12t^{\mathrm T}\Sigma_{F,N}t\bigr)$;
\item[(ii)] replacing $Q_N$ in \eqref{eq:sigma-N} by the Green pairing $Q$ of
  \eqref{eq:pairings} gives the matrix
  \begin{equation}\label{eq:sigma-F}
    \Sigma_F(j,j'):=\sqrt{w_3(k_j)w_3(k_{j'})}\;\Real\mathcal Q(\psi_j,\psi_{j'})
    =\sqrt{w_3(k_j)w_3(k_{j'})}\;\Real V\ip{(\psi_j)_L}{R(\psi_{j'})_L}_E ,
  \end{equation}
  which is a limit of the matrices \eqref{eq:sigma-N} entrywise
  (\cref{lem:hC-identification}), hence positive semidefinite as an entrywise
  limit of covariance matrices, hence the covariance matrix of a centred Gaussian
  measure $\Lambda_F$ on $\R^r$, with characteristic function
  $t\mapsto\exp\bigl(-\tfrac12t^{\mathrm T}\Sigma_Ft\bigr)$.
\end{enumerate}
Both matrices are written down here and identified in
\cref{lem:hC-identification}; this theorem asserts neither that they agree nor
that either of them is a limit.
\end{theorem}

\begin{proof}
(i) By \cref{def:gj-measure} the physical lattice law $\mu^{\mathrm{GJ}}_N$ is a
centred Gaussian measure on $\R^{I_N}$ with covariance $C_N$, and
$\omega\mapsto(\omega(\psi_1),\dots,\omega(\psi_m))$ is linear there; a linear
image of a centred Gaussian vector is a centred Gaussian vector, and by
\eqref{eq:cylinder-lift} with \eqref{eq:finite-source-law} that image is
$\mu_N^F$. For the entries: by \eqref{eq:sampled-covariance} with
\cref{lem:resolvent-pairing} the second moment
$\E_{\mu^{\mathrm{GJ}}_N}[\omega(\psi_j)\omega(\psi_{j'})]$ is
$\mathcal Q_N(\psi_j,\psi_{j'})$, a real number for real tests (\cref{sec:covariance}),
and \eqref{eq:cylinder-lift} contributes the factor
$\sqrt{w_3(k_j)w_3(k_{j'})}$; this is \eqref{eq:sigma-N}. The characteristic
function of a centred Gaussian vector with covariance $K$ is
$t\mapsto\exp(-\tfrac12t^{\mathrm T}Kt)$, the finite-dimensional Gaussian
formula, quoted and not derived. (ii) is a definition, and its second equality
is the identification $(Q=R)$ proved in \cref{thm:covariance-limit}.
\end{proof}

\begin{lemma}[The covariance entries]\label{lem:hC-identification}
In the notation of \cref{thm:char-function}:
\begin{enumerate}
\item[(i)] $\Sigma_{F,N}\to\Sigma_F$ entrywise as $N\to\infty$;
\item[(ii)] the unweighted matrix $\bigl(\Real\mathcal Q(\psi_j,\psi_{j'})\bigr)_{j,j'}$
  is the paired cosine/sine periodic-resolvent covariance matrix, and in closed
  form, for a slot pair drawn from the modes $k,l$,
  \begin{equation}\label{eq:entries}
    \begin{gathered}
      \Real\mathcal Q(\cos_k,\cos_l)=\frac{V}{2\lambda(k)}\bigl(\delta_{k,l}+\delta_{k,-l}\bigr),
      \qquad
      \Real\mathcal Q(\sin_k,\sin_l)=\frac{V}{2\lambda(k)}\bigl(\delta_{k,l}-\delta_{k,-l}\bigr),\\[2pt]
      \Real\mathcal Q(\cos_k,\sin_l)=0 ;
    \end{gathered}
  \end{equation}
\item[(iii)] for every $N$ and every slot,
  \[
    0\le\Real \mathcal Q_N(\psi_j,\psi_j)\le V/m^{2},
    \qquad\text{and}\qquad
    \Real\mathcal Q(\psi_j,\psi_j)\ge0 ,
  \]
  the second bound being an equality with $V/\bigl(2\lambda(k_j)\bigr)$ at every
  slot whose mode is non-zero, with $V/m^{2}$ at the zero-mode cosine and with
  $0$ at the zero-mode sine.
\end{enumerate}
\end{lemma}

\begin{proof}
(i) and the first half of (ii) are one computation, worth isolating because it
is where the order of the argument can be reversed.

\emph{Step 1 (the entry is a lattice second moment).} By
\eqref{eq:cylinder-lift} the coordinate $(\Pi_Fc)_j$ is
$\sqrt{w_3(k_j)}\,\omega(\psi_j)$, and $\mu^{\mathrm{GJ}}_N$ has mean zero
exactly (\cref{def:gj-measure}), so
\[
  \Sigma_{F,N}(j,j')
  =\sqrt{w_3(k_j)w_3(k_{j'})}\;\E_{\mu^{\mathrm{GJ}}_N}
     \bigl[\omega(\psi_j)\omega(\psi_{j'})\bigr],
\]
and by \eqref{eq:sampled-covariance} with \cref{lem:resolvent-pairing} the
expectation equals $\mathcal Q_N(\psi_j,\psi_{j'})$. This is \eqref{eq:sigma-N} again,
and it shows that the matrix has no source other than the lattice covariance
and the weight.

\emph{Step 2 (the limit along a fixed mode pair).} Each $\psi_j$ is a real
bounded continuous periodic test, so \cref{thm:covariance-limit} applies to
each pair: its centred-covariance line in \eqref{eq:three-limits} gives
\[
  \mathcal Q_N(\psi_j,\psi_{j'})\longrightarrow
  \Real V\ip{(\psi_j)_L}{R(\psi_{j'})_L}_E=\Real\mathcal Q(\psi_j,\psi_{j'}),
\]
the last equality being the identification $(Q=R)$ of its Step~5. Multiplying by
the fixed factor $\sqrt{w_3(k_j)w_3(k_{j'})}$ gives
$\Sigma_{F,N}\to\Sigma_F$ entrywise, which is (i) and the first half of (ii).

\emph{Step 3 (the entries, as half-sums of exponentials).} Let
$e_k(x):=e^{2\pi i\,k\cdot x/L}$, so that $\cos_k=(e_k+e_{-k})/2$ and
$\sin_k=(e_k-e_{-k})/(2i)$. The coefficient $f\mapsto F(f)$ of
\eqref{eq:fourier-coefficient} is complex linear in $f$, and at the mode $j$
the coefficient of $e_k$ is exactly $V\delta_{j,k}$: the phase is
$e^{2\pi i(k-j)\cdot x/L}$, whose integral over the physical cell vanishes
unless $k=j$ and equals $V$ when $k=j$. Inserting the half-sums into
\eqref{eq:pairings} leaves only the modes $\pm k$ and $\pm l$ of the two
coefficient vectors, and $\overline{F(e_a)}=F(e_a)$ because those coefficients
are real; summing the four surviving terms and using $\lambda(-k)=\lambda(k)$
gives the first two identities of \eqref{eq:entries}. For the mixed pair the
summand pairs the even coefficient of the cosine with the odd coefficient of
the sine, so it is odd under $p\mapsto-p$ and the total vanishes.

\emph{Step 4 (where a reader can go wrong).} \Cref{thm:covariance-limit}
already identifies the limit of precisely the pairings of Step~2, and Step~3
evaluates those pairings, so the identification is a short computation and not
a new limit theorem: no summability, tail estimate or mode-uniform envelope is
spent, because every test is a single mode. The argument can be silently
reversed here. The matrix \eqref{eq:sigma-F} must be obtained from the limit,
as in Step~2; postulating it from the symbol $\lambda^{-1}$ while invoking
\cref{thm:covariance-limit} only for the diagonal entries would be a second,
unproved assertion, and it is exactly this identification that the machine
development carries as a hypothesis on its finite-family theorem rather than
deriving from the covariance limit. The two conventions for the sign of the
imaginary part in \eqref{eq:cylinder-map} give the same matrix
\eqref{eq:sigma-F}: only the mixed entries see that sign, and
\eqref{eq:entries} makes them zero.

\emph{Step 5 (bounds).} For a real test $\psi$, Cauchy--Schwarz in
\eqref{eq:pairings}, the finite Parseval bound \eqref{eq:parseval} and the
bound $\lambda_N\ge m^{2}$ give
\[
  \mathcal Q_N(\psi,\psi)=V^{-1}\sum_{p}\frac{\abs{F_N(p)}^{2}}{\lambda_N(p)}
  \le V^{-1}\frac{V^{2}\norm{\psi}_\infty^{2}}{m^{2}}
  =\frac{V\norm{\psi}_\infty^{2}}{m^{2}},
\]
and the summands are non-negative; with $\norm{\cos_k}_\infty\le1$ and
$\norm{\sin_k}_\infty\le1$ this is the first claim of (iii). The second is
\eqref{eq:entries} read on the diagonal: it gives $V/(2\lambda(k_j))$ at every
slot whose mode is non-zero, $V/m^{2}$ at the zero-mode cosine, and $0$ at the
zero-mode sine --- the last because $\delta_{0,0}-\delta_{0,-0}=0$, which is the
right answer for a slot whose test is identically zero.
\end{proof}

\begin{theorem}[Convergence for every finite family]\label{thm:finite-family}
Let $L>0$, let $m>0$, and let $F$ be a finite family. Then
$\mu_N^F\Rightarrow\Lambda_F$ weakly on $\R^r$: for every bounded continuous
$\psi:\R^r\to\R$, $\int\psi\,\dd\mu_N^F\to\int\psi\,\dd\Lambda_F$ as
$N\to\infty$.
\end{theorem}

\begin{proof}
Three steps, none of which mentions the carrier.

\emph{(i) The matrices converge.} This is \cref{lem:hC-identification}(i).

\emph{(ii) The family is tight in $\R^r$, uniformly in $N$.} By
\cref{lem:hC-identification}(iii) each diagonal entry satisfies
$0\le\Sigma_{F,N}(j,j)\le V/m^{2}=L^{2}/m^{2}$, so the second moment
$\E_{\mu_N^F}\abs{x}^{2}=\operatorname{tr}\Sigma_{F,N}$ is at most
$rL^{2}/m^{2}$ for every $N$. Markov's inequality gives
$\mu_N^F\{x:\abs{x}>R\}\le rL^{2}/(m^{2}R^{2})\to0$ as $R\to\infty$, and
closed balls in $\R^r$ are compact, so the family is tight. A single Gaussian
law on $\R^r$ needs no estimate for this; what the bound supplies is tightness
of the whole family at once.

\emph{(iii) The characteristic functions converge.} By
\cref{thm:char-function} the characteristic function of $\mu_N^F$ is
$t\mapsto\exp(-\tfrac12t^{\mathrm T}\Sigma_{F,N}t)$, which by (i) converges for
each fixed $t$ to $\exp(-\tfrac12t^{\mathrm T}\Sigma_Ft)$, that of $\Lambda_F$.
The limit is continuous at the origin and is the characteristic function of a
probability measure, so Levy's continuity theorem --- pointwise convergence of
characteristic functions to a function continuous at the origin implies weak
convergence --- gives $\mu_N^F\Rightarrow\Lambda_F$. This is the standard fact
that centred Gaussian laws on $\R^r$ with convergent covariance matrices
converge weakly \cite{bogachev1998}, and the boundedness of the variances in
(ii) is its licence: uniformly bounded diagonal entries give tightness, and
\cref{lem:hC-identification}(iii) says the limit does not degenerate.
\end{proof}

What \cref{thm:finite-family} establishes is convergence of the pushforwards
under every finite coordinate map; by \cref{def:cylinder-map} these are the laws
of the paired cosine/sine coordinates, with the resolvent matrix as limit. What
it does not establish is anything about measures on the carrier $\mathcal C$. A
family of finite-dimensional limits does not by itself produce a measure on
$\mathcal C$, and convergence of every marginal would not imply weak
convergence on $\mathcal C$ without a tightness hypothesis; hence the reliance
on the uniform moment and compact tails of \cref{sec:tightness} and on the
identification of the two sigma algebras in \cref{thm:cylinder-borel}. The
passage runs: tightness supplies a limit point (\cref{thm:tightness}), continuity
of $\Pi_F$ transfers \cref{thm:finite-family} to that limit point, and
\cref{thm:cylinder-uniqueness} says two Borel probability measures with the same
finite marginals agree. Nothing here produces a measure on the carrier.

\begin{remark}[The index dictionary, once]\label{rem:index-dictionary}
The family $F$ of \cref{def:cylinder-map} does not depend on $N$: its modes are
integers, fixed once and for all, and only the mesh varies. This is where the
two indexing conventions of \cref{hyp:volume} meet, and the dictionary is stated
here once, with $N=n+1$ and $a=a_n=L/(n+1)$: the law
\eqref{eq:finite-source-law} at index $n$ is $\mu_{n+1}$, and
\cref{thm:finite-family} is the statement that these laws converge as
$n\to\infty$. Taking the finite family \emph{before} the mesh limit is what
makes \eqref{eq:sigma-F} mesh-free, and it is why the theorem concerns a
sequence of laws on one and the same space $\R^r$. By \cref{hyp:volume} the
spacing is a function of the index and not a free parameter, so no double limit
occurs in this section.
\end{remark}
